\mxnewpage\section*{Correction EV3-T1}\label{EV3-T1-CORR}
\textbf{QCM :} 1 A, 2 C, 3 D, 4 D, 5 C, 6 C, 7 C, 8 B, 9 B, 10 C, 11 C, 12 A. Chaque réponse correcte : \num{0,5} point ; total 6.
\paragraph{Vérifications des automatismes.} Priorités ; somme de fractions ; soustraction des exposants ; coefficient multiplicateur \num{1,2} ; division par \num{0,8} ; produit \num{1,25} fois \num{0,8} ; résolution affine ; accroissement vertical divisé par l’horizontal ; moyenne ; angle divisé par 360 ; complément à 1 ; substitution entre parenthèses.
\subsection*{Exercice 1 -- Un gain composé}
\begin{enumerate}
\item Valeurs -2,1,4 ; probabilités $\dfrac{1}{4}$,$\dfrac{1}{2}$,$\dfrac{1}{4}$. \textit{Barème : 1 point. Résultat et justification élémentaire attendus ; demi-point possible pour démarche correcte inachevée.}
\item $E(X)=1$ ; $E(X^2)=1+\dfrac{1}{2}+4=\dfrac{11}{2}$ ; $V(X)=\dfrac{9}{2}$. \textit{Barème : 2 points. Méthode ou justification : 1 ; calcul et conclusion : 1.}
\item Gain brut moyen 3 : mise équitable 3. Écart type inchangé $\dfrac{3}{\sqrt 2}$. \textit{Barème : 2 points. Méthode ou justification : 1 ; calcul et conclusion : 1.}
\end{enumerate}
\subsection*{Exercice 2 -- Pic d’un modèle}
\begin{enumerate}
\item $f\prime=-(t+1)e^{-t}$. \textit{Barème : 1 point. Résultat et justification élémentaire attendus ; demi-point possible pour démarche correcte inachevée.}
\item Strictement décroissante sur $\intfo{0}{+\infty}$, maximum 2 en 0. \textit{Barème : 2 points. Méthode ou justification : 1 ; calcul et conclusion : 1.}
\item $y=2-t$ ; approximation \num{1,9}. Ce n’est pas une égalité exacte. \textit{Barème : 2 points. Méthode ou justification : 1 ; calcul et conclusion : 1.}
\end{enumerate}
\subsection*{Exercice 3 -- Cercle trigonométrique}
\begin{enumerate}
\item $A\pt{\dfrac{\sqrt 3}{2};\dfrac{1}{2}}$, $B\pt{-\dfrac{\sqrt 3}{2};\dfrac{1}{2}}$. \textit{Barème : 1 point. Résultat et justification élémentaire attendus ; demi-point possible pour démarche correcte inachevée.}
\item $AB=\sqrt3$. \textit{Barème : 1 point. Résultat et justification élémentaire attendus ; demi-point possible pour démarche correcte inachevée.}
\item Milieu (0;$\dfrac{1}{2}$), AB horizontal ; médiatrice verticale $x=0$. \textit{Barème : 2 points. Méthode ou justification : 1 ; calcul et conclusion : 1.}
\end{enumerate}
\paragraph{Étalonnage.} Durée cible : 20–25 min pour les automatismes, 80–85 min pour les exercices, 10–15 min de relecture. Ces durées sont des estimations éditoriales à mesurer en classe. Accepter toute méthode correcte ; ne pas sanctionner deux fois une erreur numérique propagée avec cohérence. Le barème proposé ne vaut pas décision d’harmonisation de l’équipe.
